\uuid{DbQZ}
\titre{DL en l'infini/Calcul de limites/Asymptotes}
\auteur{hanani ; mzari}
\organisation{exo7}
\licence{CC BY-NC-SA 4.0}
\datecreate{2020-01-20}
\niveau{L1}
\module{Analyse}
\chapitre{Développement limité}
\sousChapitre{Applications}
\difficulte{3}
\theme{Applications des DL, DL en l'infini, Calcul de limites, Asymptotes}
\exo7qcmid{lille-2.developpements-limites.047}

\contenu{
\texte{
Soit $a\in \Rr$ et $f$ une fonction. On dit que $f(x)$ est un $o\left((x-a)^n\right)$ au voisinage de $a$ et on écrit $f(x)=o((x-a)^n)$ si $\displaystyle \lim _{x\to a}\frac{f(x)}{(x-a)^n}=0$. On note par $DL_n(a)f(x)$ le développement limité de $f$ à l'ordre $n$ en $a$. On note aussi $DL_n(a^+)f(x)$ (resp. $DL_n(a^-)f(x)$) le développement limité de $f$ à l'ordre $n$ à droite en $a$ (resp. à gauche en $a$).
}

\begin{question}
\qclassification{125.03}
\qtags{motcle=DL en l'infini/Calcul de limites/Asymptotes}
On considère les fonctions $f$ et $g$ telles que $\displaystyle f(x)=\sqrt{x^2+1}\mathrm{e}^{1/(x+1)}$ et $\displaystyle g(x)=xf\left(\frac{1}{x}\right)$. On note $\Gamma$ le graphe de $f$ et $\Delta$ la droite d'équation $y=x+1$. Parmi les affirmations suivantes, cocher celles qui sont vraies :
\begin{answers}
\good{$\displaystyle DL_3(0^+)g(x)=1+x+\frac{2}{3}x^3+o(x^3)$.}
\bad{$\displaystyle \lim _{x\to 0}g(x)=1$ et $g$ se prolonge par continuité en $0$.}
\good{Au voisinage de $+\infty$, $\Gamma$ admet la droite $\Delta$ comme asymptote et $\Gamma$ est situé au dessus de $\Delta$.}
\bad{Au voisinage de $-\infty$, $\Gamma$ admet la droite $\Delta$ comme asymptote et $\Gamma$ est situé au dessus de $\Delta$.}
\end{answers}
\vskip2mm
\begin{explanations}
On a : $\displaystyle g(x)=\sqrt{1+x^2}\mathrm{e}^{x/(1+x)}$ si $x>0$ et $\displaystyle g(x)=-\sqrt{1+x^2}\mathrm{e}^{x/(1+x)}$ si $x<0$. On calcule le $DL_3(0)\sqrt{1+x^2}\mathrm{e}^{x/(1+x)}$. Ce qui donne :
$$DL_3(0^+)g(x)=1+x+\frac{2}{3}x^3+o(x^3)\Rightarrow DL_2(+\infty)f(x)=x+1+\frac{2}{3x^2}+o\left(\frac{1}{x^2}\right)$$
et
$$DL_3(0^-)g(x)=-1-x-\frac{2}{3}x^3+o(x^2)\Rightarrow DL_2(-\infty)f(x)=-x-1-\frac{2}{3x^2}+o\left(\frac{1}{x^2}\right).$$
Ainsi, au voisinage de $+\infty$, $\Gamma$ admet $\Delta$ comme asymptote et $\Gamma$ est situé au dessus de $\Delta$ car $\displaystyle f(x)-y=\frac{2}{3x^2}+o\left(\frac{1}{x^2}\right)\geq 0$. Mais, au voisinage de $-\infty$, $\Gamma$ admet la droite $\Delta '$ d'équation $y=-x-1$ comme asymptote, et $\Gamma$ est situé en dessous de $\Delta '$ car $\displaystyle f(x)-y=-\frac{2}{3x^2}+o\left(\frac{1}{x^2}\right)\leq 0$. Enfin, $g$ n'admet pas de limite en $0$ car $\displaystyle \lim _{x\to 0^+}g(x)=1$ et $\displaystyle \lim _{x\to 0^-}g(x)=-1$.
\end{explanations}
\end{question}
}
